\begin{chapintro}
組み入れた試験の論文から、統合に使う数字と試験の特徴を取り出し、それぞれの試験がどれくらい偏りにくいかを評価します。この章では、抽出フォームの作り方、本文の記載から 2×2 の数を作る方法、RoB 2 によるバイアスのリスクの評価を説明し、取り出したデータを R に読み込むところまで進めます。
\end{chapintro}

\section{何を取り出すか}
\label{sec:04-1}

データの抽出は、試験ごとに同じ項目を同じ形で書き写す作業です。項目は、あらかじめ抽出フォーム（表計算ソフトの表で十分です）にしておきます。フォームを先に作るのは、試験ごとに取り出す項目がぶれないようにするためで、検査でいえば検査依頼書の様式を統一しておくのと同じ役割を果たします。降圧薬 A のレビューでは、次の項目を取り出します。

\begin{description}
\item[試験の識別] 筆頭著者、発表年、試験名、試験登録番号を書き、同じ試験の論文が複数あれば一覧にします。
\item[試験の方法] 研究デザイン、実施国と施設数、追跡期間、ランダム化と割り付けの方法を書きます。
\item[対象] 組み入れ基準、人数、年齢、性別、試験開始時の血圧、脳卒中の既往がある人の割合を書きます。
\item[介入と比較] 降圧薬 A の用量、プラセボの内容、併用を認めた薬剤を書きます。
\item[アウトカム] 脳卒中の定義と判定方法、評価の時点、群ごとの発症数と解析対象の人数を書きます。
\item[バイアスのリスク] RoB 2 の各領域の判定と、その根拠になった本文の記載を書きます（この章の後半で説明します）。
\end{description}

抽出も、選別と同じく 2 人が独立に行い、食い違いを話し合って解決します。数字の書き写しは単純に見えますが、群の取り違え、表の行の読み違い、時点の取り違えが意外に多く起こります。2 人で独立に抽出して突き合わせれば、こうした誤りの多くを見つけられます。転記した数字は、どの論文の何ページの何の表から取ったかも記録しておきます。

\section{本文から 2×2 の数を作る}
\label{sec:04-2}

脳卒中のような「起きた、起きなかった」で数えるアウトカムでは、群ごとに発症した人数と解析した人数の 2 つずつ、合わせて 4 つの数を取り出します。この 4 つの数を並べたものを 2×2 表と呼びます（表 4.1）。

\begin{kkcenter}
\kkfit{\begin{tabular}{lccc}
\toprule
 & 脳卒中あり & 脳卒中なし & 合計 \\
\midrule
降圧薬 A 群 & 121 & 1,479 & 1,600 \\
プラセボ群 & 186 & 1,414 & 1,600 \\
\bottomrule
\end{tabular}}
\kkcap{表 4.1　架空の Garcia 2018 の 2×2 表}
\end{kkcenter}

いちばん楽なのは、本文に人数がそのまま書かれている場合です。架空の Garcia 2018 の結果の段落に「脳卒中は降圧薬 A 群 1,600 人中 121 人（7.6\%）、プラセボ群 1,600 人中 186 人（11.6\%）に発生した」とあれば、121、1,600、186、1,600 をそのまま写します。括弧の中の割合は、121 ÷ 1,600 = 0.0756 から計算できるので、転記の誤りを確かめるのに使えます。

割合しか書かれていない場合は、割合と人数から発症数を逆算します。架空の Fischer 2016 の本文に「脳卒中の発症は降圧薬 A 群 800 人の 4.0\%、プラセボ群 800 人の 4.8\%」とあれば、降圧薬 A 群は 800 × 0.040 = 32 人、プラセボ群は 800 × 0.048 = 38.4 人と計算できます。人数は整数なので 38 人とします。割合が丸められているときは、逆算した値が整数からずれることがあります。そのときは、ずれの大きさから元の人数を絞り込み、本文の別の箇所（表や図の注）とも照らし合わせて確かめます。

どの人数を分母にするかにも注意が要ります。RCT の結果は、割り付けた全員をそのまま割り付けた群で解析する方法（ITT 解析、intention-to-treat）と、計画どおりに薬剤を飲み続けた人だけで解析する方法（per protocol 解析）の 2 通りで報告されることがあります。途中で薬剤をやめた人を除くと、やめた理由（副作用や病状の悪化）が群によって違う場合に、ランダム化でそろえた背景が崩れます。効果の推定には ITT 解析の数を使うのが原則です。

結果がグラフ（累積発症曲線など）にしか示されていない場合は、まず著者に問い合わせます。それでも数が得られなければ、グラフから値を読み取るソフトを使うこともありますが、読み取った値であることを記録し、感度分析で除外したときの結果も確かめます（\kkref{sec:07-6}）。

\section{バイアスのリスクを評価する}
\label{sec:04-3}

RCT はランダム化によって両群の背景をそろえますが、試験の進め方によっては結果が偏ります。それぞれの試験がどれくらい偏りやすいかを評価する道具として、RCT では RoB 2（Sterne 2019）が使われています。RoB 2 は、結果を偏らせる経路を 5 つの領域に分け、領域ごとにバイアスのリスクを「低い」「懸念あり」「高い」の 3 段階で判定します。

\begin{description}
\item[ランダム化の過程] 割り付けの順番が予測できない方法で作られ、その順番が登録の担当者から隠されていたか（割り付けの隠蔽）。降圧薬 A の試験で、担当者が次の患者はプラセボ群だと前もって分かるようでは、重症の患者を降圧薬 A 群へ割り当てるといった操作ができてしまいます。隠蔽は、割り付ける時点で順番を伏せることです。
\item[意図した介入からの逸脱] 割り付けたあとに、参加者や治療者がどちらの群かを知っていたために、治療の内容が変わっていないか。降圧薬 A の錠剤がプラセボと見た目で区別できると、医師が降圧薬 A 群の患者にだけ他の降圧治療を控える、といったことが起こりえます。どちらの群かを伏せておくことを盲検化と呼び、隠蔽とは時点が違います。
\item[アウトカムデータの欠測] 脱落や追跡不能で結果が分からない人が多すぎないか、その理由が群で違わないか。降圧薬 A 群だけ副作用で来院をやめた人が多ければ、残った人の脳卒中の数は比べられません。
\item[アウトカムの測定] 脳卒中かどうかを判定した人が、患者がどちらの群かを知らずに判定したか。知っていると、判定の境界の症例で判断がぶれます。
\item[報告される結果の選択] 計画段階で決めたアウトカムと解析を、そのとおりに報告しているか。血圧の変化も脳卒中も測っていた試験が、結果のよかったほうだけを報告していないかを見ます。
\end{description}

5 つの領域に分けて点検するやり方は、誤報告の原因を検査前、検査、検査後の工程に分けて調べるのに似ています。試験全体の判定は、5 つの領域の判定から決めます。どれか一つでも「高い」があれば全体も高い、すべてが「低い」なら全体も低い、それ以外は「懸念あり」とするのが基本の規則です。ただし、複数の領域で懸念ありが重なり、結果への信頼を大きく損なうと判断すれば、全体を「高い」とすることもできます。

降圧薬 A の 10 試験の判定を図 4.1 に示します。架空の判定ですが、理由の書き方を 3 試験ぶん示します。架空の Aoki 2009 は、割り付けの順番の作り方は書かれているものの、隠蔽の方法の記載がないので、ランダム化の過程が「懸念あり」です。架空の Chen 2012 は、脳卒中の判定者が群を知らされていたかどうかの記載がなく、アウトカムの測定が「懸念あり」です。架空の Endo 2015 は、隠蔽の記載がないうえ、割り付けた 180 人のうち 38 人が途中で追跡できなくなっていました。論文はその人たちを脳卒中なしとみなして、分母を 90 人ずつのまま解析しており、追跡できなかった理由も群ごとに示されていません。数の流れは、割り付け 180 人、追跡できず 38 人、追跡できて結果が分かる 142 人です。表 1.1 の 6 人と 12 人は、結果が分かる 142 人から数えた脳卒中で、追跡できなかった 38 人は 0 人として分母にだけ入っています。追跡できなかった人に脳卒中が起きていれば結果が変わるので、アウトカムデータの欠測を「高い」と判定し、試験全体も「高い」としました。

\begin{kkcenter}
\begin{minipage}{\linewidth}\centering
\begin{kkfitpic}
\begin{tikzpicture}[x=10.5mm, y=5.2mm]
\foreach \lab [count=\j] in {D1,D2,D3,D4,D5,全体}
  \node[font=\scriptsize\bfseries] at (\j+1.6, 0) {\lab};
\foreach \name/\a/\b/\c/\d/\e/\o [count=\i] in {
  Aoki 2009/m/l/l/l/l/m, Berger 2011/l/l/l/l/l/l, Chen 2012/l/l/l/m/l/m,
  Dubois 2014/l/l/l/l/l/l, Endo 2015/m/l/h/l/l/h, Fischer 2016/l/l/l/l/l/l,
  Garcia 2018/l/l/l/l/l/l, Hayashi 2019/l/l/l/l/m/m, Ivanov 2021/l/l/l/l/l/l,
  Jensen 2022/l/l/l/l/l/l} {
  \node[anchor=west, font=\scriptsize] at (-0.9, -\i) {\name};
  \foreach \v [count=\j] in {\a,\b,\c,\d,\e,\o} {
    \if\v l \fill[kksecondary] (\j+1.6,-\i) circle (1.8mm);
      \node[font=\tiny] at (\j+1.6,-\i) {+};\fi
    \if\v m \fill[kkpale] (\j+1.6,-\i) circle (1.8mm);
      \draw[kkaccent] (\j+1.6,-\i) circle (1.8mm);
      \node[font=\tiny] at (\j+1.6,-\i) {?};\fi
    \if\v h \fill[kkwarm] (\j+1.6,-\i) circle (1.8mm);
      \node[font=\tiny, text=white] at (\j+1.6,-\i) {−};\fi
  }
}
\end{tikzpicture}
\end{kkfitpic}
\kkcap{図 4.1　架空の 10 試験の RoB 2 の判定。D1 ランダム化の過程、D2 意図した介入からの逸脱、D3 アウトカムデータの欠測、D4 アウトカムの測定、D5 報告される結果の選択。＋ 低い、？ 懸念あり、− 高い}
\end{minipage}
\end{kkcenter}

判定には、必ず本文の該当箇所を根拠として書き添えます。「記載がない」ことも根拠の一つです。論文に書かれていない手順は、実際に行われていても評価者には確かめようがないので、記載がなければ懸念ありとするのがふつうです。必要なら著者や試験登録の記録で確かめます。

バイアスのリスクの判定は、点数にして合計しないでください。5 つの領域は結果を偏らせる経路が違い、どの領域の問題がどれだけ結果を動かすかも試験ごとに違うからです。評価の結果は、\kkref{sec:06-4} のサブグループ解析と \kkref{sec:07-6} の感度分析で、判定の違う試験の間で結果が変わるかを調べるのに使います。

\begin{hattenbox}
この冊子の例は RCT だけですが、介入のレビューでは非ランダム化研究を含めることもあります。その場合の評価には ROBINS-I（Sterne 2016）を使います。交絡や対象者の選び方による偏りを、仮想的なランダム化試験と比べる形で評価する道具です。改訂版の ROBINS-I V2 は、riskofbias.info で 2025 年 11 月版が草案として公開されていますが、確定版の論文は確認できませんでした。使うときは最新の版と手引きを確かめてください。
\end{hattenbox}

\section{抽出したデータを R に読み込む}
\label{sec:04-4}

抽出した数は、1 行を 1 試験とし、列に項目を並べた表にして、CSV 形式（項目をカンマで区切ったテキスト）で保存します。列名は英字で短く付けておくと、R で扱いやすくなります。降圧薬 A のデータでは、試験名（study）、対象（pop）、バイアスのリスクの全体判定（rob）、降圧薬 A 群の発症数と人数（ev\_t、n\_t）、プラセボ群の発症数と人数（ev\_c、n\_c）の 7 列にしました。

R で CSV を読み込むには read.csv() を使います。日本語を含むファイルは、文字コードを指定して読み込みます。

\begin{rcode}
bp <- read.csv("data/bp_stroke.csv",
               fileEncoding = "UTF-8")
str(bp)    # 列の名前と型を確かめる
\end{rcode}

\begin{rout}
'data.frame':	10 obs. of  7 variables:
 $ study: chr  "Aoki 2009" "Berger 2011" "Chen 2012" ...
 $ pop  : chr  "高齢者" "一般" "高齢者" "一般" ...
 $ rob  : chr  "懸念あり" "低" "懸念あり" "低" ...
 $ ev_t : int  7 48 44 134 6 32 121 10 112 43
 $ n_t  : int  150 1200 420 2500 90 800 1600 300 3000 650
 $ ev_c : int  17 65 39 145 12 38 186 23 139 83
 $ n_c  : int  150 1200 420 2500 90 800 1600 300 3000 650
\end{rout}

str() の出力で、10 行 7 列のデータが読み込まれ、発症数と人数が整数（int）として入っていることを確かめます。数が文字（chr）になっていたら、どこかの値に数字以外の文字（全角数字や空白）が紛れ込んでいます。読み込んだら、発症数が人数を超えていないか、合計人数（10 試験で 21,420 人）が抽出フォームと一致するかも確かめておきます。

\section{この章のまとめ}
\label{sec:04-5}

データは、あらかじめ作った抽出フォームに 2 人が独立に書き写し、出典の箇所も記録します。発症数は ITT 解析の数を使い、割合しかなければ人数から逆算します。各試験のバイアスのリスクは RoB 2 の 5 つの領域で判定し、根拠を書き添え、点数にはしません。抽出したデータは 1 行 1 試験の CSV にして R に読み込みました。次の章では、この 10 試験のリスク比を R で統合します。
